\documentclass[conference]{IEEEtran}

\usepackage{cite}
\usepackage{amsmath,amssymb,amsfonts}
\usepackage{graphicx}
\usepackage{textcomp}
\usepackage{parskip}
\usepackage{xcolor}

\begin{document}

\title{Active Learning in Neural Networks: A Comparison of Learning Approaches}

\author{
\IEEEauthorblockN{Aphelele M. Zenzile}
\IEEEauthorblockA{
Stellenbosch, South Africa \\
27403122@sun.ac.za
}
}

\maketitle

\begin{abstract}
    The effectiveness of active learning approaches relative to conventional passive learning remains 
    dependent on the learning strategy and problem being considered. This study compares passive learning 
    using stochastic gradient descent with two active learning approaches, neural network output sensitivity analysis 
    and uncertainty sampling, under a consistent experimental framework. The three approaches were implemented and 
    evaluated using common performance criteria across three classification and three function approximation problems 
    of varying complexity. The experimental results were analysed to compare the performance of the learning approaches 
    under the specified conditions. This paper therefore presents a comparative study of passive and active learning 
    in neural networks, examining how the three approaches perform across problems with different learning characteristics 
    and identifying the observed advantages and limitations of active learning relative to passive learning.
\end{abstract}

\begin{IEEEkeywords}
active learning, neural networks, stochastic gradient descent,
uncertainty sampling, sensitivity analysis
\end{IEEEkeywords}


\section{Introduction}

The effectiveness of neural network training depends not only on the model architecture and optimisation procedure, 
but also on how training data are selected and presented to the learning algorithm. Conventional passive learning approaches 
expose a model to the available training data without actively considering which observations are most useful for improving the model. 
Active learning instead seeks to make the training process more selective by identifying observations that are expected to provide 
greater value for learning. The distinction between passive and active learning has therefore become relevant when considering 
how neural networks can achieve effective learning while making more deliberate use of available training data.

Different active learning strategies use different criteria for determining which examples should be selected. 
A comparison is therefore required to determine how these strategies perform relative to passive learning and 
whether the use of active selection provides measurable performance advantages across problems of different complexity.

This study addresses the problem by implementing a common underlying neural network training algorithm on top of which three learning strategies 
are incorporated: passive learning using stochastic gradient descent (SGD), active learning using neural network output sensitivity analysis, 
and active learning using uncertainty sampling. The three strategies are evaluated through a consistent set of experiments and common 
performance criteria, allowing their performance to be compared under the same experimental framework. The implementation used for the experiments 
is available in this public repository: \textit{https://github.com/APZenzile/ML\_Assignment\_3}. Three classification problems and 
three function approximation problems are used to evaluate the approaches, with the selected problems representing varying levels of complexity. 
This provides a basis for examining whether the relative performance of the learning strategies is consistent across different learning tasks.

The remainder of the report is organised as follows. Section~\ref{sec:background} presents the background concepts underlying passive and active 
learning and the learning strategies considered in the study. Section~\ref{sec:methodology} describes the experimental methodology, including the 
selected problems, data preprocessing, neural network architecture, regularisation, training approaches, control parameter selection, and performance 
evaluation procedure. Section~\ref{sec:results} presents the experimental results and comparison of the three learning strategies. Section~\ref{sec:conclusion} 
concludes the report by summarising the main findings of the comparative study and acknowledging the limitations and contraints under which the study was conducted.

\section{Background}
\label{sec:background}


% Background content


\section{Methodology}
\label{sec:methodology}

% Methodology content

\subsection{Problem Description}

% Classification and function approximation problems


\subsection{Data Preprocessing}

% Preprocessing procedures and motivation


\subsection{Neural Network Architecture}

% Network architecture, hidden units, activation and cost functions


\subsection{Regularisation}

% Weight decay and penalty coefficient selection


\subsection{Training Approaches}

% Passive learning using SGD
% Active learning using sensitivity analysis
% Active learning using uncertainty sampling


\subsection{Control Parameter Selection}

% Selection of algorithm control parameters


\subsection{Performance Evaluation}

% Performance criteria and empirical statistical comparison


\section{Results}
\label{sec:results}


% Experimental results


\section{Conclusion}
\label{sec:conclusion}


% Conclusion


\bibliographystyle{IEEEtran}
\bibliography{references}

\end{document}